%%%%%%%%%%%%%%%%%%%%%%%%%%%%%%%%%%%%%%%%%

%%%%%%%%%%%%%%%%%%%%%%%%%%%%%%%%%%%%%%%%%

%----------------------------------------------------------------------------------------
%	PACKAGES AND OTHER DOCUMENT CONFIGURATIONS
%----------------------------------------------------------------------------------------

\documentclass[final]{beamer}
\usepackage{SAMBaPoster}
\usepackage{PosterText}
\usepackage{../../Documents/Documents_Academic/Reports/notation}

%----------------------------------------------------------------------------------------
%	TITLE SECTION 
%----------------------------------------------------------------------------------------
\columnsep=20pt % This is the amount of white space between the columns in the poster
\columnseprule=5pt

%%%%%%%%%%%%%%%%%%%%%%%EDITEDIT
% Fill in the following fields for the title section. 
% To format the display of the title section you will need to edit the style file beamerconfposter.sty.

%----------------------------------------------------------------------------------------
\usebackgroundtemplate{\includegraphics[height=\paperheight]{wifi3}}
\begin{document}

%----------------------------------------------------------------------------------------
%%%%%%%%%%%%%%%%%%%%%%%EDITEDIT 
% To change the background image\usebackgroundtemplate{\includegraphics[width=\paperwidth, height=\paperheight]{INSERT IMAGE NAME}
%\usebackgroundtemplate{\includegraphics[width=\paperwidth]{wifi}}
\begin{frame}[t] % The whole poster is enclosed in one beamer frame

\begin{columns}[t] % The whole poster consists of three major columns, the second of which is split into two columns twice - the [t] option aligns each column's content to the top

%\begin{column}{\sepwid}\end{column} % Empty spacer column

\begin{column}{\firstcolwid} % The first column

%---------------------------------------------------------------
% Introduce the Method
%---------------------------------------------------------------
\intro

%------------------------------------------------------------------------------
%	Explain the method
%--------------------------------------------------------------------------------

\MathematicalMotivation
\application
%----------------------------------------------------------------------------------------
%	Introduce The Project
%----------------------------------------------------------------------------------------

%---------------------------------------------------------------

\end{column} % End of the first column

%\begin{column}{1.0\sepwid}\end{column} % Empty spacer column

\begin{column}{\twocolwid} % Begin a column which is two columns wide (column 2)
%----------------------------------------------------------------------------------------

  % Coloured straight line box
% \vspace{-2cm}
% \hspace{0.5in}
% \begin{beamercolorbox}[wd=14in,colsep=0.15cm]{cboxb}
% \end{beamercolorbox}

\trajectories
\methodblock
\end{column}
%----------------------------------------------------------------
%		End the first column and start the second
%------------------------------------------------------------

%\begin{column}{\sepwid}% SPACER COLUMN\end{column}

\begin{column}{\threecolwid}
\coverageevaluation

\compareblock

%\futurework
%----------------------------------------------------------------------------------------
%----------------------------------------------------------------------------------------
%	REFERENCES
%----------------------------------------------------------------------------------------
\biblio

%\begin{column}{\sepwid}
%% SPACER COLUMN
%\end{column}
%----------------------------------------------------------------------------------------
%	ACKNOWLEDGEMENTS
%----------------------------------------------------------------------------------------

% \begin{block}{Acknowledgements}

% \end{block}

%---------------------------------------------------------------

\end{column} % End of the third column

\end{columns} % End of all the columns in the poster
\end{frame} % End of the enclosing frame

\end{document}
